\documentclass[instructions.tex]{subfiles}
\begin{document}
 
\chapter{De l’humilité,}
 
\primo Peu de gens réfléchissent sur la nécessité de cette vertu, et peu de gens la pratiquent. Vertu néanmoins si nécessaire, que sans elle on ne se sauve point. Elle est, selon saint Bernard, le fondement et la gardienne de toutes les vertus. Si elle ne les précède, dit saint Augustin, si elle ne les accompagne, si elle ne les suit, l’orgueil les emporte. Vertu si essentielle, que le Fils de Dieu est descendu du ciel pour l’enseigner aux hommes : Apprenez de moi, dit-il, non à faire des miracles, mais apprenez de moi que je suis doux et humble de cœur.
 
Voilà la première leçon de ce grand maître. Il en a donné lui-même l’exemple ; il s’est anéanti jusqu’à prendre la forme d’esclave ; il s’est confondu parmi les pécheurs et les ignorans. Quoiqu’il fût la lumière du monde et le Saint des saints, il n’a jamais eu de complaisance en lui-même\footnote{Christus non sibi placuit. Rom. 15. 3.}.
 
Il a passé sa vie dans l’obscurité, et n’a paru en public que pour réparer par ses humiliations la gloire de son Père, que l’orgueil de l’homme avait flétrie. Il a choisi pour sa mère la plus pure des vierges ; mais elle n’aurait jamais été digne de cet honneur, si elle n’avait pas été la plus humble des créatures, et Jésus-Christ aurait mieux aimé avoir pour mère une femme du commun, que de naître d’une superbe. O humilité d’un Dieu, que tu confonds notre orgueil !
 
Il en est de l’humilité comme de la foi ; sans elle on ne peut plaire à Dieu, ni entrer dans le ciel. Non, dit Jésus-Christ, si vous ne devenez comme de petits enfans, vous n’entrerez pas dans le royaume des cieux. Voyez l’humilité d’un enfant, il ne s’estime pas plus que les autres : ses petits talens, la noblesse, les richesses, ne lui enflent pas le cœur ; il ne se vante point ; si on le loue, il en rougit, il sent son ignorance et sa foiblesse ; aimant à s’instruire, il ne raisonne et ne chicane point sur ce qu’on lui dit ; simple, soumis, il reconnoît ses fautes, et souffre qu’on le reprenne. Un enfant est sincère, il n’est ni défiant ni soupçonneux ; sans ressentiment, sans prévention, il ne parle mal de personne, et respecte chacun. Voilà le modèle que le Fils de Dieu nous propose. Ce sont les ames simples et dociles avec lesquelles il aime à s’entretenir\footnote{Cum simplicibus sermocinatio ejus. Prov. 3. 32.}; tandis qu’il résiste aux superbes, leur refuse ses grâces, les abat comme des arbres morts et infructueux. Les pluies et les rosées ne s’arrêtent pas sur les hauteurs, mais dans la profondeur des vallées ; ainsi les grâces du ciel ne s’arrêtent point dans les cœurs superbes, mais dans les cœurs profondément humbles.
 
\secundo Il n’est personne qui n’aime l’humilité dans les autres, mais il en est peu qui la pratiquent. Pour s’attirer l’estime, la bienveillance du monde, on sait cacher son ignorance, déguiser ou excuser des défauts qu’on a, pour faire paroître des talens et des vertus qu’on n’a pas. On n’ose pas toujours se vanter, mais on mendie des louanges, en se blâmant soi-même, dans la vue d’être excusé ou applaudi. Orgueil hypocrite !
 
Cachez vos défauts avec prudence, vous le devez ; l’humilité ne consiste pas à les faire paroître. Cachez-les, non pas pour être estimé, mais afin de ne pas scandaliser. L’humilité ne consiste pas aussi à cacher toutes ses vertus et ses talens. Qu’on voie vos bonnes œuvres, dit Jésus-Christ, qu’on sache que vous vivez chrétiennement, que vous employez utilement vos talens et vos biens, non pas afin qu’on vous loue, mais afin qu’on glorifie votre Père céleste.
 
Celui qui est vraiment humble reconnoît que tout ce qu’il a de talens et de richesses vient de Dieu. Il reconnoît aussi de bonne foi ce qu’il y a de défectueux et de mauvais en lui-même ; il en gémit, il souffre qu’on l’avertisse et qu’on le reprenne ; il ne souffre qu’avec peine qu’on l’estime et qu’on le loue : semblable à un malade qui reconnoît qu’il est malade, et ne peut souffrir qu’on dise qu’il est en santé ; et de même qu’un malade dit toujours son mal plus grand, et prend même plaisir à l’entendre dire ; aussi l’homme parfaitement humble se fait un plaisir d’entendre ce qui l’humilie.
 
C’est une chose étrange, qu’il ne soit personne qui craigne plus le mépris que celui qui en est le plus digne, et que les plus vicieux, les plus décriés ne veuillent point être repris, et soient les plus délicats sur le point d’honneur. Un vrai chrétien a des sentimens bien différens : il ne craint point la confusion, et ne croit jamais qu’on lui fait tort ; parce qu’il croit mériter plus de mal qu’on ne lui en fait et qu’on ne lui en dit ; c’est pourquoi chacun est plein d’estime pour lui, parce qu’il a de l’estime pour tous, excepté pour lui-même.
 
\secundo La vraie humilité est fondée sur la connoissance de Dieu et sur la connoissance de soi-même. Voilà la science des saints sur la terre, science la plus importante. De quoi sert la connoissance du mouvement des astres et des secrets de la nature, si l’on se méconnoît soi-même, si l’on ignore ce que l’on est devant Dieu, et ce que l’on a mérité ?
 
Quand on connoît les grandeurs et la sainteté de Dieu, et qu’on connoît sa propre bassesse et les misères de son ame, oh ! qu’on se voit misérable, et qu’on comprend combien on est peu de chose ! L’estime et les mépris des hommes ne touchent plus une ame qui est ainsi pénétrée de sa bassesse et de son néant; c’est alors qu’on est véritablement grand devant Dieu. Voyez l’humble publicain : loin de s’applaudir comme le pharisien, le poids de ses crimes l’accable de confusion ; il s’accuse, il se fait des reproches, il s’humilie; et Dieu l’élève, lui accorde sa miséricorde, tandis que l’orgueilleux pharisien est rejeté de Dieu.
 
\section{Résolutions}
 
\primo Je demanderai souvent l’humilité à Dieu.

\secundo Je lui rapporterai la gloire de tout le bien qui est en moi, et de tout ce que je ferai de bon.

\tertio Je parlerai le moins que je pourrai de moi et de ce qui me regarde, et ce ne sera jamais ni pour me vanter, ni pour être loué.

\quarto Je supporterai avec charité et résignation les humiliations, pensant que je les ai méritées par mes péchés.
 
\prayer{O mon Dieu ! donnez-moi l’humilité, et faites que cette vertu préside désormais à tous mes discours, à tout ce que je ferai et souffrirai.}
 
\end{document}
